The rapid growth of digital education has highlighted the urgent need for scalable,
intuitive, and pedagogically effective tools for creating interactive learning content.
At Vietnam National University -- Ho Chi Minh City (VNU-HCM), the primary approach for
producing H5P interactive videos relies on the WordPress plugin ecosystem, a solution
that imposes significant technical barriers, suffers from performance limitations, and
results in adoption rates as low as 23\% among faculty.

This thesis presents the design, development, and evaluation of a purpose-built,
cloud-native H5P interactive video content creation platform tailored for university
instructors. The platform integrates a React~18 and TypeScript frontend, a
Node.js/Express backend, a SQLite relational database managed via the Sequelize ORM,
and a dual-provider artificial intelligence pipeline that combines Anthropic Claude
with a locally hosted Ollama inference server. The AI pipeline analyses video
transcripts---sourced from YouTube automatic captions, uploaded caption files, or
audio transcription via OpenAI Whisper---and generates contextually appropriate H5P
interactive elements including Multiple Choice, True/False, Fill in the Blanks,
Hotspot, Drag-and-Drop, and Mark the Words questions, all classified according to
Bloom's revised taxonomy. The system applies structured prompt engineering and
runtime schema validation via Zod to guarantee that every AI-generated suggestion
conforms to the H5P content type specification before it is presented to the
instructor for review.

A comparative usability evaluation conducted with ten Information Technology students
acting as instructor proxies demonstrated substantial improvements over the
WordPress-based baseline: task completion rate increased from 30\% to 95\%, median
content creation time fell from 45~minutes to 3.2~minutes, and user satisfaction
scores rose from 2.1 to 4.8 on a five-point scale. Security controls were aligned to
the OWASP Top Ten through JWT-based authentication, bcrypt password hashing
(cost factor~12), HTTP security headers via Helmet, strict CORS allowlisting, and
input validation at every API boundary.

The thesis contributes a reproducible architecture for AI-assisted interactive
content authoring, an empirically validated workflow that reduces cognitive overhead
for non-technical educators, and a practical integration path for the VNU-HCM MOOC
infrastructure. The findings demonstrate that user-centred design combined with
generative AI can fundamentally lower the barrier to interactive pedagogy in
Southeast Asian higher education.

\bigskip
\noindent
\textbf{Keywords:} H5P, interactive video, educational technology, AI-assisted
content authoring, user-centred design, usability evaluation, LMS integration,
prompt engineering, Vietnamese higher education.
